\documentclass[12pt,a4paper]{article}
\usepackage[utf8]{inputenc}
\usepackage[T1]{fontenc}
\usepackage[spanish]{babel}
\usepackage{geometry}
\geometry{a4paper, top=2.5cm, bottom=2.5cm, left=2.5cm, right=2.5cm}
\usepackage{listings}
\usepackage{xcolor}
\usepackage{tcolorbox}
\usepackage{hyperref}
\usepackage{graphicx}
\usepackage{tikz}
\usetikzlibrary{positioning, arrows.meta}

\definecolor{codegreen}{rgb}{0,0.6,0}
\definecolor{codegray}{rgb}{0.5,0.5,0.5}
\definecolor{codepurple}{rgb}{0.58,0,0.82}
\definecolor{backcolour}{rgb}{0.96,0.96,0.96}

\lstdefinestyle{javastyle}{
    language=Java,
    backgroundcolor=\color{backcolour},
    commentstyle=\color{codegreen}\itshape,
    keywordstyle=\color{blue}\bfseries,
    numberstyle=\tiny\color{codegray},
    stringstyle=\color{codepurple},
    basicstyle=\ttfamily\footnotesize,
    breakatwhitespace=false,
    breaklines=true,
    captionpos=t,
    keepspaces=true,
    numbers=left,
    numbersep=5pt,
    showspaces=false,
    showstringspaces=false,
    showtabs=false,
    tabsize=4,
    frame=single,
    rulecolor=\color{lightgray}
}
\lstset{
    style=javastyle,
    literate={á}{{\'a}}1 {é}{{\'e}}1 {í}{{\'i}}1 {ó}{{\'o}}1 {ú}{{\'u}}1
             {Á}{{\'A}}1 {É}{{\'E}}1 {Í}{{\'I}}1 {Ó}{{\'O}}1 {Ú}{{\'U}}1
             {ñ}{{\~n}}1 {Ñ}{{\~N}}1
}

\title{\textbf{Guía de Solución: Programación Orientada a Objetos} \\ \large Actividad 4: Manejo de Excepciones y Archivos}
\author{Cristian}
\date{}

\begin{document}
\maketitle

\section{Repositorio de GitHub}
El código fuente de este proyecto, integrando las Actividades 3, 4 y 5 mediante un menú gráfico unificado, se encuentra alojado en GitHub. 
\newline
\textbf{Enlace del repositorio:} \url{https://github.com/cioranemil/poo-trabajo-3-}

\section{Enunciado: Actividad 4}
Se requiere desarrollar una aplicación gráfica interactiva que contenga y dé solución a 5 ejercicios específicos planteados en el libro guía sobre el manejo de excepciones y la lectura de flujos de datos (archivos). El programa debe asegurar que cualquier error en la lógica, entrada de datos o gestión de archivos sea controlado elegantemente sin interrumpir la ejecución, empleando bloques \texttt{try/catch/finally} y directivas \texttt{throw}.

\section{Casos de Uso}
\begin{table}[h]
    \centering
    \begin{tabular}{|p{3cm}|p{12cm}|}
    \hline
    \textbf{Caso de Uso} & \textbf{Gestión de Excepciones Interactivas} \\ \hline
    \textbf{Actor} & Usuario \\ \hline
    \textbf{Descripción} & El usuario interactúa con la interfaz para forzar diversos errores lógicos o de validación, observando cómo el sistema reacciona y controla la eventualidad. \\ \hline
    \textbf{Flujo Principal} & 
    1. El usuario inicia la aplicación y selecciona ``Actividad 4 - Excepciones''.\newline
    2. En el menú, elige uno de los ejercicios (Ej: Probar try/catch o Validar Vendedor).\newline
    3. El usuario ingresa datos intencionalmente erróneos (ej. letras donde van números o edad menor a 18).\newline
    4. El sistema ejecuta el bloque \texttt{try}.\newline
    5. El sistema detecta la falla, salta al bloque \texttt{catch} o lanza un \texttt{throw}, e informa al usuario del problema. \\ \hline
    \textbf{Flujo Alternativo} & 
    Si el usuario ingresa datos correctos (ej. Vendedor mayor de 18 años), la aplicación confirma el éxito de la operación demostrando el funcionamiento estable del programa. \\ \hline
    \end{tabular}
    \caption{Descripción del Caso de Uso Principal (Actividad 4)}
\end{table}

\section{Diagrama de Clases y Explicación}
\begin{figure}[h]
    \centering
    \begin{tikzpicture}[
        node distance=2.5cm,
        umlbox/.style={draw, inner sep=0pt, align=left, font=\footnotesize}
    ]
    
    \node[umlbox] (menu) {
        \begin{tabular}{l}
        \multicolumn{1}{c}{\textbf{VentanaPrincipalActividad4}} \\
        \hline
        - btnExcepciones: JButton \\
        - btnVendedor: JButton \\
        - btnMatematicas: JButton \\
        \hline
        <<constructor>> VentanaPrincipalActividad4() \\
        - abrir(JFrame ventana)
        \end{tabular}
    };
    
    \node[umlbox, below left=1.5cm and -1cm of menu] (excepcion) {
        \begin{tabular}{l}
        \multicolumn{1}{c}{\textbf{VentanaExcepciones}} \\
        \hline
        - areaTexto: JTextArea \\
        \hline
        + ejecutarSimulacion() \\
        \end{tabular}
    };
    
    \node[umlbox, below right=1.5cm and -1cm of menu] (vendedor) {
        \begin{tabular}{l}
        \multicolumn{1}{c}{\textbf{Vendedor}} \\
        \hline
        - nombre: String \\
        - apellidos: String \\
        - edad: int \\
        \hline
        <<constructor>> Vendedor(...) \textit{throws Exception}
        \end{tabular}
    };
    
    \draw[-stealth] (menu) -- node[left, near end] {Abre} (excepcion);
    \draw[-stealth] (menu) -- node[right, near end] {Crea obj} (vendedor);
    
    \end{tikzpicture}
    \caption{Diagrama de clases parcial de la Actividad 4}
\end{figure}

\textbf{Explicación del diagrama de clases:}
Se presenta un diagrama simplificado que ilustra el patrón modular usado. La clase \texttt{VentanaPrincipalActividad4} funge como despachador. Está asociada con múltiples submódulos como \texttt{VentanaExcepciones} (para probar bloques try/catch mediante una simulación de consola visual en un \texttt{JTextArea}) y a clases lógicas independientes como \texttt{Vendedor}, cuyo constructor está diseñado específicamente para lanzar excepciones (\texttt{IllegalArgumentException}) si no se cumplen las reglas de negocio (ser mayor de 18 años).

\section{Diagrama de Objetos}
\begin{verbatim}
[ menuAct4: VentanaPrincipalActividad4 ] ------------ [ uiExcepcion: VentanaExcepciones ]
\end{verbatim}

\section{Ejecución del Programa}
A continuación se muestra la interfaz principal del programa durante su ejecución:

\begin{figure}[h]
    \centering
    % \includegraphics[width=0.6\textwidth]{interfaz_act4.png}
    \vspace{4cm} 
    \caption{Interfaz gráfica de la Actividad 4 probando excepciones (reemplace la imagen)}
\end{figure}

\section{Solución: Código de Muestra (Gestión de Errores)}
Se adjunta el extracto más importante de la validación del constructor de la clase \texttt{Vendedor}, ilustrando el lanzamiento intencional de excepciones para cumplir el objetivo del ejercicio 6.5.

\begin{lstlisting}[caption={Clase Vendedor.java}]
package Vendedor;

public class Vendedor {
    private String nombre;
    private String apellidos;
    private int edad;

    public Vendedor(String nombre, String apellidos, int edad) {
        this.nombre = nombre;
        this.apellidos = apellidos;
        
        // Objetivo: Generar y lanzar excepcion controlada si la regla falla
        if (edad < 18) {
            throw new IllegalArgumentException("El vendedor debe ser mayor de edad.");
        }
        
        this.edad = edad;
    }
}
\end{lstlisting}

\end{document}
